% Notas de aula — Capítulo 16: Ciclones Tropicais
% Curso: Meteorologia Prática (baseado em R. Stull, Practical Meteorology, CC BY-NC-SA 4.0)
% Caixas verdes = conteúdo literal do quadro; linhas verdes = encenação.
\documentclass[11pt]{article}
\usepackage{../comum/preambulo}
\graphicspath{{../figuras/cap16/}}

\title{Capítulo 16 --- Ciclones Tropicais\\
{\large Notas de aula (roteiro de quadro-negro)}}
\author{Curso de Meteorologia Prática}
\date{}

\begin{document}
\maketitle
\thispagestyle{fancy}

\noindent\textbf{Plano sugerido:} 2 aulas de 2\,h.
\emph{Aula 1:} \S\ref{sec:anatomia}--\S\ref{sec:carnot} (anatomia,
condições de gênese, a máquina de Carnot).
\emph{Aula 2:} \S\ref{sec:dinamica}--\S\ref{sec:perigos} (dinâmica
radial, trajetórias, maré de tempestade e o caso Catarina).

\section{Um motor diferente de tudo que vimos}\label{sec:anatomia}

\begin{noquadro}[abertura --- contraste com o Cap.~13]
\textbf{Cap.~16 --- Ciclones Tropicais}\\[2pt]
\begin{tabular}{lll}
 & extratropical (Cap.~13) & tropical (hoje)\\
energia & baroclinia ($\Delta T$ horizontal) & calor latente do
oceano\\
núcleo & frio & \textbf{quente}\\
frentes & sim & não\\
vento máx.\ & altitude (jato) & \textbf{superfície} (parede do olho)\\
\end{tabular}
\end{noquadro}

Anatomia: olho de 20--50\,km (subsidência, céu claro, calmaria),
\textbf{parede do olho} (o anel de Cb com os ventos máximos), bandas
espirais de chuva, e o escoamento anticiclônico de saída na tropopausa
\javimos{5}{a tampa}. Tudo axissimétrico — sem as frentes do mundo
extratropical \javimos{12}{frentes}.

\begin{noquadro}[condições de gênese (a lista clássica)]
1. TSM $\ge 26{,}5\,°$C numa camada de $\ge$50\,m\\
2. distância do equador $\gtrsim 5°$ (precisa de $f$!
\javimos{10}{Coriolis})\\
3. cisalhamento vertical FRACO ($< 10$\,m/s) — o contrário do Cap.~14!\\
4. umidade na média troposfera $+$ perturbação inicial (onda de leste)
\end{noquadro}

Essa lista explica o mapa das trajetórias — e por que o Atlântico Sul
quase não tem furacões (cisalhamento alto $+$ TSM marginal): o
Catarina (2004) foi a exceção instrutiva.

\section{A máquina de Carnot}\label{sec:carnot}

\quadro{Desenhar o ciclo no corte do furacão: entrada isotérmica junto
ao mar (pega calor), subida adiabática úmida na parede, saída
isotérmica no topo frio, descida na periferia. É um diagrama de
Carnot desenhado no céu.}

\begin{formadif}[intensidade potencial (T1; Casos~1 e 3)]
Motor entre $T_s$ (mar) e $T_{out}$ (topo da bigorna):
\[ v_{max}^2 = \frac{C_k}{C_d}\,
\frac{T_s - T_{out}}{T_{out}}\,\Delta k, \]
com $\Delta k$ o déficit de entalpia ar--mar. O fator
$(T_s{-}T_{out})/T_{out} \simeq 0{,}5$ supera o Carnot clássico porque
o calor dissipado pelo atrito é reciclado. O limiar dos 26{,}5\,°C é
Clausius--Clapeyron \javimos{4}{7\%/K}: $\Delta k$ explode com a TSM.
\end{formadif}

\begin{noquadro}[WISHE --- a realimentação]
vento forte $\to$ mais evaporação (fórmula bulk!
\javimos{3}{fluxos bulk}) $\to$ mais calor latente na parede $\to$
núcleo mais quente $\to$ pressão mais baixa
\javimos{1}{hipsométrica} $\to$ vento mais forte\\[2pt]
teto: a intensidade potencial (MPI) — calculada de verdade com o
\texttt{tcpyPI} no Caso~3
\end{noquadro}

O núcleo quente sustenta a baixa hidrostaticamente:
$\Delta T = 10$\,K em 10\,km $\Rightarrow$ $\Delta P \sim 40$\,hPa
(T3). Ciclones tropicais são \emph{os} sistemas de núcleo quente; os
do Cap.~13, de núcleo frio — a distinção aparece limpa nos cortes de
satélite \javimos{8}{sondagem por micro-ondas}.

\section{Dinâmica radial: Holland e o vento gradiente}
\label{sec:dinamica}

\begin{formadif}[perfil de Holland (T2; Caso~2)]
\[ P(r) = P_c + \Delta P\,e^{-(R_w/r)^B},\qquad
v(r) = \sqrt{\frac{r}{\rho}\frac{dP}{dr} + \Big(\frac{fr}{2}\Big)^2}
- \frac{|f|r}{2}. \]
Vento máximo na parede ($r \simeq R_w$); dentro do olho o vento
desaba — é o vento gradiente do Cap.~10 \javimos{10}{vento gradiente}
levado ao limite ($V^2/R$ dominando).
\end{formadif}

Trajetórias: alísios empurram para oeste \javimos{11}{alísios}; a
deriva beta (os \emph{beta gyres} — parentes das ondas de Rossby
\javimos{11}{$\beta$}) soma 1--3\,m/s para oeste e para o polo; nas
latitudes médias os oestes capturam e o sistema recurva (T5). Ao
tocar terra: sem oceano, o motor desliga em $\sim$1 dia — mas a chuva
continua.

\section{Perigos: a água mata mais que o vento}\label{sec:perigos}

\quadro{Ranking no quadro, perguntando antes à turma qual mata mais:
1º maré de tempestade, 2º chuva/enchentes, 3º vento, 4º tornados
embutidos.}

\begin{formadif}[maré de tempestade (T4; Caso~4)]
Vento empilha o mar contra a costa; no equilíbrio:
\[ \frac{d\eta}{dx} = \frac{\tau}{\rho_{agua}\,g\,h(x)},\qquad
\tau = \rho_{ar}C_D V^2, \]
mais o barômetro invertido ($\sim$1\,cm/hPa — T4). Plataforma rasa e
larga amplifica: 50\,m/s dão $\sim$8\,m de surge em plataforma rasa
(N4) — a física por trás de Katrina e Bangladesh.
\end{formadif}

Escala Saffir--Simpson (CT1--CT5 por vento máximo): lembrar que dano
$\propto v^3$ (potência do vento) e a chuva não escala com a
categoria — sistemas lentos e fracos podem ser os mais mortais
(Harvey). Num clima $+2$\,K, a MPI sobe poucos \%/K, mas chuva
($+7$\%/K) e nível do mar compõem o risco \veremos{21}{ciclones no
clima futuro} (N3).

\section{Síntese da aula}

\begin{noquadro}[mapa final --- canto do quadro]
\[ v_{max}^2 = \frac{C_k}{C_d}\frac{T_s - T_{out}}{T_{out}}\Delta k
\qquad P(r) = P_c + \Delta P e^{-(R_w/r)^B} \]
\[ \frac{d\eta}{dx} = \frac{\tau}{\rho g h} \qquad
\text{gênese: TSM} \ge 26{,}5\,°\mathrm{C},\ |\phi| \gtrsim 5°,\
\text{cisalhamento fraco} \]
núcleo quente, motor de Carnot, morte em terra — e a água mata mais
que o vento
\end{noquadro}

\section{Exercícios complementares}\label{sec:exercicios}

Soluções (professor) em \texttt{cap16\_solucoes.ipynb}.

\subsection*{Teóricos}

\begin{enumerate}[label=\textbf{T\arabic*.},itemsep=4pt]
  \item Compare a eficiência de Carnot clássica $(T_s-T_{out})/T_s$
        com o fator de Emanuel $(T_s-T_{out})/T_{out}$ para
        $T_s = 300$\,K, $T_{out} = 200$\,K, e explique fisicamente a
        diferença (reciclagem do calor dissipado).
  \item Do perfil de Holland, obtenha o vento ciclostrófico
        $v^2 = (r/\rho)\,dP/dr$ e mostre que o máximo fica em
        $r \approx R_w$.
  \item Use a equação hipsométrica para estimar a queda de pressão à
        superfície produzida por um núcleo 10\,K mais quente numa
        camada de 10\,km.
  \item Mostre que o barômetro invertido dá
        $d\eta/dP = 1/(\rho_{agua}g) \approx 1$\,cm/hPa e calcule a
        contribuição para um CT5 ($\Delta P = 90$\,hPa).
  \item Explique a deriva beta: por que um vórtice intenso num planeta
        com $\beta$ migra para oeste e para o polo mesmo sem
        escoamento médio?
\end{enumerate}

\subsection*{Numéricos (no notebook)}

\begin{enumerate}[label=\textbf{N\arabic*.},itemsep=4pt]
  \item Mapeie a MPI no plano (TSM, $T_{out}$); quanto vale esfriar o
        topo em 5\,K vs.\ aquecer o mar em 1\,K?
  \item Ajuste ($P_c$, $R_w$, $B$) de Holland para vento máximo de
        60\,m/s em 25\,km e classifique o furacão.
  \item Com o \texttt{tcpyPI}, aqueça sondagem e TSM em $+2$\,K e
        calcule a intensificação em \%/K.
  \item Varra largura e profundidade da plataforma no modelo de surge
        e mapeie a elevação na costa para $V = 50$\,m/s.
\end{enumerate}

\vspace{1em}
\noindent\rule{\linewidth}{0.4pt}\\
{\scriptsize Material derivado de R.~Stull, \emph{Practical Meteorology}
(CC BY-NC-SA 4.0); este documento herda a mesma licença.}

\end{document}
